%% kapa: exam=1 round=25 year=2014 period=1 form=A subject=civil_law subject_ko=민법 seq=10 number=10
%% source: Q-Net 감정평가사 기출문제 (cst00309) artlSeq=5101207 file=extracted/1008/5101207/2014년 제25회 감정평가사 1차시험 기출문제/1교시(A형).hwp
%% conversion: kordoc 4.24.51 -> tools/gen1.py (auto); review=proofread-sonnet; flags=none
\begin{question}{10}
등기의 유효요건에 관한 설명으로 옳지 않은 것은? (다툼이 있으면 판례에 의함)
\begin{choices}
  \item 중복된 소유권보존등기의 등기명의인이 동일인이 아닌 경우, 선등기가 원인무효가 아닌 한 후등기는 무효이다.
  \item 무효등기의 유용에 관한 추인은 언제나 명시적으로 이루어져야 한다.
  \item 무효인 중복등기에 터 잡은 등기부시효취득은 인정되지 않는다.
  \item 중간생략등기의 사법상 효력이 모든 경우에 부정되는 것은 아니다.
  \item 소유권에 관한 등기가 적법한 원인 없이 말소되더라도 원칙적으로 그 소유권의 효력에는 아무런 영향을 미치지 않는다.
\end{choices}
\end{question}
\answer{2}
